\section{Finite-horizon sample--memory bounds}
\label{sec:v22-frontier}
The minimum $W_c$ in Theorem~\ref{thm:v21-endtoend} permits arbitrarily
long finite training. To retain the preparation constraint, let
$\mathsf V_{N,\boldsymbol K}(B,Q)$ denote the supremum of the worst-case
expected score over the same reset auditors, using at most $N$ training
trials and satisfying the specified register bound at each acquisition,
random-bit, decision and validation cut. The observation order and the
external clock are fixed parts of this interface. Write
$\mathsf V_{N,W}$ when every bound equals $W$. Padding an audit with
ignored trials shows monotonicity in $N$. This value does not identify
machines having only the same terminal dictionary.

A second realization of a response cover gives a polynomial preparation
bound. Its counter is larger than the four-state block counter, and
both constructions remain useful. The elementary random-walk argument
below is classical; the purpose of retaining its proof is to price its
use in the complete, calibrated continuation machine.

\begin{lemma}[A truncated absorbing comparison]\label{lem:v22-walk}
Let $Z_i$ be independent Bernoulli variables with mean $u$. Fix integers
$m,n\ge1$ and $0<\theta<1/2$. Start at $S_0=0$, add $2Z_i-1$
while $|S_{i-1}|<m$, and retain the state forever after either boundary
is reached. At time $n$ declare positive exactly when $S_n\ge0$.
The machine has $2m+1$ states. Whenever $|u-1/2|\ge\theta$, its error
probability is at most
\begin{equation}\label{eq:v22-walk-error}
 e_{m,n}(\theta)=e^{-4\theta m}+e^{-2n\theta^2}.
\end{equation}
If $\tau$ is the boundary hitting time without truncation, then, with
$x=2u-1$,
\begin{equation}\label{eq:v22-walk-time}
 \E_u\tau=
 \begin{cases}
 m\tanh(m\operatorname{arctanh}x)/x,&0<|x|<1,\\
 m^2,&x=0,\\
 m,&|x|=1.
 \end{cases}
\end{equation}
In particular $\E_u(\tau\wedge n)\le\min(n,m^2)$, and outside the
indifference band it is at most $\min(n,m^2,m/(2\theta))$.
\end{lemma}
\begin{proof}
First suppose $1/2<u<1$ and put $r=(1-u)/u$. For the unabsorbed walk
$T_j=\sum_{i\le j}(2Z_i-1)$, the process $r^{T_j}$ is a nonnegative
martingale. Stopping on $\{-m,m\}$ gives
\[
 \Pp_u(T_\tau=-m)=\frac{r^m}{1+r^m}\le r^m.
\]
For completeness, absorption has finite mean: in each block of $2m$
steps, a run of $2m$ positive increments has probability $u^{2m}>0$
and forces absorption from every interior state. Thus the unabsorbed
survival time is bounded by a geometric number of such blocks. Stopping
first at $\tau\wedge j$, then using bounded convergence on the stopped
exponential martingale, justifies the displayed identity. The estimate
$\log((1+2\theta)/(1-2\theta))\ge4\theta$ gives
$r^m\le e^{-4\theta m}$ when $u\ge1/2+\theta$.

If the walk has not been absorbed, a negative final answer requires
$T_n<0$. Hoeffding's inequality~\cite{Hoeffding1963} for variables in $[-1,1]$ gives
\[
 \Pp_u(T_n\le0)\le e^{-2n(u-1/2)^2}\le e^{-2n\theta^2}.
\]
Partitioning a wrong answer according to absorption proves
\eqref{eq:v22-walk-error}. Reflection proves the lower-mean case; its
nonabsorbed error includes $T_n=0$ and satisfies the same bound. At
$u=0,1$ the assertion follows directly.

The drift martingale $T_j-xj$ stopped at $\tau\wedge j$ gives
$\E_u T_\tau=x\E_u\tau$, by integrability of $\tau$ and boundedness
of the stopped walk. The hitting probabilities yield
$\E_u T_\tau=m(1-r^m)/(1+r^m)$. Since
$r=\exp(-2\operatorname{arctanh}x)$, this is \eqref{eq:v22-walk-time}
when $x\ne0$. For $x=0$, stop $T_j^2-j$ to obtain $\E\tau=m^2$.
The endpoint formula is deterministic. Finally
$\tanh(my)\le m\tanh y$ for integer $m\ge1$ and $y\ge0$:
induction follows from the addition formula for $\tanh$. This proves
$\E\tau\le m^2$; the bound $m/|x|$ follows from $|T_\tau|\le m$.
Truncation can only decrease the mean.
\end{proof}

\begin{theorem}[A preparation--width family of realizations]
\label{thm:v22-tournament}
Let $B\subset\Delta(\Omega)$ be compact, $d=|\Omega|$, and let
$h_1,\ldots,h_K\in[0,1]^\Omega$, $K\ge2$, have envelope
$v=\min_{b\in B}\max_i(Q-b)h_i$. Round their coordinates to dyadic
precision $r$, obtaining $h_i^{(r)}$, and choose known dyadic numbers
$q_i\in[0,1]$ with $|q_i-Qh_i^{(r)}|\le\rho$. Suppose $q_i$ also have
denominator $2^r$; enlarging $r$ accommodates any specified dyadic
calibration. For $m,n\ge1$ and $0<\theta<1/2$, an explicitly clocked
fair-bit audit satisfies
\begin{align}
 N&=n(K-1),\qquad W\le3dK(2m+1),\label{eq:v22-tournament-resources}\\
 \inf_{b\in B}\E_bD
 &\ge v-2^{1-r}-2\rho-4(K-1)\theta
                 -2(K-1)e_{m,n}(\theta).\label{eq:v22-tournament-score}
\end{align}
At completed training-trial cuts its width is at most $K(2m+1)$;
at the decision cut it is at most $K$. Every random-bit cut, the
selected event, and both validations are included in $W$. In particular
one can replace $e_{m,n}(\theta)$ by $\delta$ by taking
\begin{equation}\label{eq:v22-polynomial}
 m=\left\lceil\frac{\log(2/\delta)}{4\theta}\right\rceil,
 \qquad
 n=\left\lceil\frac{\log(2/\delta)}{2\theta^2}\right\rceil.
\end{equation}
\end{theorem}
\begin{proof}
For a fixed unknown $b$, set
$g_i(b)=q_i-bh_i^{(r)}$. Compare the current incumbent $i$ with the next
index $j$ on fresh observations. Conditional on the observation $w$,
generate a coin with probability
\begin{equation}\label{eq:v22-comparison-coin}
 \frac12+\frac{q_j-q_i-h_j^{(r)}(w)+h_i^{(r)}(w)}4.
\end{equation}
This is in $[0,1]$, is dyadic with denominator $2^{r+2}$, and has mean
$1/2+(g_j(b)-g_i(b))/4$. Apply Lemma~\ref{lem:v22-walk}, selecting
$j$ for a positive answer and $i$ otherwise. Conditional on all earlier
stages, the observations and fair bits for this comparison are fresh.
The probability of selecting the worse index when
$|g_j-g_i|\ge4\theta$ is at most $e_{m,n}(\theta)$.

Outside these exceptional comparisons each tournament stage loses at
most $4\theta$ relative to the better of its two contestants. Induction
shows a total loss at most $4(K-1)\theta$ relative to
$\max_i g_i$. The union probability of exceptional comparisons is at
most $(K-1)e_{m,n}(\theta)$ even though the incumbent is random.
All $g_i$ lie in $[-1,1]$, so their range is at most two. Consequently
\[
 \E_b g_{\widehat i}(b)\ge\max_i g_i(b)-4(K-1)\theta
                                -2(K-1)e_{m,n}(\theta).
\]
The calibration changes selection and evaluation by at most $2\rho$.
For each coordinate rounding, $\|h_i-h_i^{(r)}\|_\infty\le2^{-r}$,
so $|(Q-b)(h_i-h_i^{(r)})|\le2^{1-r}$. A nested random event with mean
$h_{\widehat i}^{(r)}$ has exactly that rounded payoff; hence
\eqref{eq:v22-tournament-score}.

At a completed trial retain the incumbent and the absorbing walk state.
The stage number $j$ is supplied by the predetermined clock, not by a
sample-dependent register. During coin generation additionally retain
$w$ and at most three sampler labels, by
Proposition~\ref{prop:v21-rational}. This gives $3dK(2m+1)$ labels.
The terminal event sampler, held event, and validation storage use at
most $K(2d+1)$, $K(d+1)$, and $2K(d+1)$ labels, respectively, as in
Theorem~\ref{thm:v21-endtoend}; all are bounded by the stated $W$.
The stage boundary discards the walk and initializes the next one.
After absorption, remaining observations and bits in that stage are
ignored, but are still charged to the fixed schedule. Thus the count
$N$ includes padding. Each active coin uses $r+2$ fair bits. A coarse
bound for the full number of epochs is $T\le N(r+3)+r+4$.

The finite transition description consists of the dyadic coin table,
the boundary and horizon integers, and the nested-event thresholds. It
has $O(K^2d(r+3)+\log m+\log n+\log K)$ bits under binary encoding;
the constants depend on a fixed encoding convention. No unknown $b$
occurs in that description. The autonomous conversion stores the clock
and has width at most $(T+1)W$, rather than $W$. Substitution of
\eqref{eq:v22-polynomial} proves the last assertion.
\end{proof}

\begin{corollary}[Two general preparation regimes]
\label{cor:v22-two-realizations}
For a strict $K$-response cover and a prescribed positive score margin,
there are whole-schedule rational realizations of that margin with
width at most $12dK$ and exponential block preparation count, and with
width at most $3dK(2m+1)$ and preparation count $n(K-1)$ as in
\eqref{eq:v22-polynomial}. Both constrain every training cut. Their
resource pairs are upper bounds for $\mathsf V_{N,W}$, not a claim of
an optimal frontier for every compact image.
\end{corollary}
\begin{proof}
Choose rounding, calibration, band and error budgets whose sum is
smaller than the strict margin. Apply Theorem~\ref{thm:v21-endtoend}
or Theorem~\ref{thm:v22-tournament}. A one-response cover needs no
training and is implemented directly by the rational event sampler.
\end{proof}
